\subsection{正项级数的收敛判别法}

在本节中，我们把（滥用语言的说法）这样的实项级数 \((u_{n})\) ：其项 \(u_{n} \geqslant 0\) 从 \(n\) 的某个特定值起成立，理解为正项级数。我们关于这类级数所说的一切，通过改变符号，立即可以推广到各项都 \(\leqslant 0\) （从 \(n\) 的某个特定值起）的级数。我们已经看到 (II, p.~64, 例 3)，对点的每个序列 \((\mathbf{u}_n)_{n\geqslant 1}\) （取值于赋范空间 \(E\) ），可以关联一个阶梯函数 \(\mathbf{u}\) ，它定义在 \([1, +\infty[\) 上，由条件 \(\mathbf{u}(x) = \mathbf{u}_n\) （对 \(n\leqslant x < n + 1\) ）确定：这时级数 \((\mathbf{u}_n)\) 收敛当且仅当积分 \(\int_1^{+\infty}\mathbf{u}(t)dt\) 收敛。

设 \((u_{n})\) 与 \((v_{n})\) 是两个正项级数，u 与 v 是相应的阶梯函数：关系 \(u_{n} \leqslant v_{n}\) （对 \(n \geqslant n_{0}\) ）等价于 \(u(x) \leqslant v(x)\) （对 \(x \geqslant n_{0}\) ）。于是关系 \(u_{n} \preccurlyeq v_{n}, u_{n} \prec v_{n}, u_{n} \sim v_{n}\) 中的每一个分别等价于 \(u(x) \preccurlyeq v(x), u(x) \prec v(x)\) 与 \(u(x) \sim v(x)\) ；这一注记使我们能把命题 1 (V, p.~228) 与命题 6 (V, p.~230) 表述如下：

命题 1. 设 \((u_{n})\) 与 \((v_{n})\) 是两个正项级数。若 \(u_{n} \preccurlyeq v_{n}\) 且级数 \((v_{n})\) 收敛，则 \((u_{n})\) 收敛；若 \(u_{n} \succcurlyeq v_{n}\) 且 \(\sum_{n=1}^{\infty} v_{n} = +\infty\) ，则 \(\sum_{n=1}^{\infty} u_{n} = +\infty\) 。

命题 2. 设 \((u_{n})\) 与 \((v_{n})\) 是两个正项级数：

\(1^{\circ}\) 若级数 \((v_{n})\) 收敛，则关系 \(u_{n} \ll v_{n}\) （相应地， \(u_{n} \sim v_{n}\) ）蕴含 \(\sum_{p=n}^{\infty} u_{p} \ll \sum_{p=n}^{\infty} v_{p}\) （相应地， \(\sum_{p=n}^{\infty} u_{p} \sim \sum_{p=n}^{\infty} v_{p}\) ）。

\(2^{\circ}\) 若 \(\sum_{n=1}^{\infty} v_n = +\infty\) ，则关系 \(u_n \ll v_n\) （相应地， \(u_n \sim v_n\) ）蕴含 \(\sum_{p=1}^{n} u_p \ll \sum_{p=1}^{n} v_p\) （相应地， \(\sum_{p=1}^{n} u_p \sim \sum_{p=1}^{n} v_p\) ）。

在命题 1 中取比较级数 \((v_{n})\) 为通项形如 \(v_{n} = f(n)\) 的级数（其中 \(f\) 是一个 \(\geqslant 0\) 的函数，对每个实数 \(x > x_0\) 有定义且在区间 \([x_0, +\infty[\) 上递减），便得到方便的收敛判别法；事实上：

命题 3（柯西–麦克劳林判别法）. 若 \(f\) 是一个 \(\geqslant 0\) 且在 \([x_0, +\infty[\) 上递减的函数，则以 \(v_n = f(n)\) 为通项的级数收敛当且仅当积分 \(\int_{x_0}^{+\infty} f(t) dt\) 收敛。

为证明这一点，只需注意：若 v 是与级数 \((v_{n})\) 相关联的阶梯函数，则由于 f 递减， \(v(x+1) \leqslant f(x) \leqslant v(x)\) 对每个 \(x \geqslant x_{0}\) 成立；于是该命题是比较原理 (II, p.~66, prop. 3) 的推论。

由于出现在积分的对数收敛判别法 (V, p.~229, prop. 2、3 与 4) 中的函数在某个区间 \([x_{0}, +\infty[\) 上递减，应用 V, p.~237 的命题 1 与命题 3，就得到下列判别法：

命题 4（“0 阶对数判别法”）. 设 \((u_{n})\) 是正项级数；若 \(u_{n} \preccurlyeq n^{\mu}\) 对某个 \(\mu < -1\) 成立，则级数 \((u_{n})\) 收敛；若 \(u_{n} \succcurlyeq n^{\mu}\) 对某个 \(\mu \geqslant -1\) 成立，则级数 \((u_{n})\) 的和为无穷。

命题 5（“p 阶对数判别法”）. 设 \((u_{n})\) 是正项级数。若 \(u_{n} \preccurlyeq \frac{1}{nl_{1}(n)l_{2}(n)\ldots l_{p-1}(n)(l_{p}(n))^{\mu}}\) 对某个 \(\mu > 1\) 成立，则级数 \((u_{n})\) 收敛；若 \(u_{n} \succcurlyeq \frac{1}{nl_{1}(n)l_{2}(n)\ldots l_{p-1}(n)(l_{p}(n))^{\mu}}\) 对某个 \(\mu \leqslant 1\) 成立，则级数 \((u_{n})\) 的和为无穷。

若 \(0 \leqslant q < 1\) ，则 \(q^n \preccurlyeq n^{-\mu}\) 对任何 \(\mu > 0\) 成立；再应用 0 阶对数判别法便证明了几何级数 \(\sum_{n=0}^{\infty} q^n\) 对 \(|q| < 1\) 收敛 (Gen.~Top., IV, p.~364)。若在命题 1 中取 \(v_n = q^n\) ，就得到一个可以写成如下形式的判别法（“柯西判别法”）：设 \((u_n)\) 是正项级数；若 \(\limsup_{n \to \infty} (u_n)^{1/n} < 1\) ，则级数 \((u_n)\) 收敛；若 \(\limsup_{n \to \infty} (u_n)^{1/n} > 1\) ，则级数 \((u_n)\) 的和为无穷。事实上，若 \(\limsup_{n \to \infty} (u_n)^{1/n} = a < 1\) ，则 \(u_n \preccurlyeq q^n\) 对每个 \(q\) （使 \(a < q < 1\) ）成立。反之，若 \(\limsup_{n \to \infty} (u_n)^{1/n} = a > 1\) ，则 \(u_n \geqslant q^n > 1\) 对无穷多个 \(n\) 的值成立（对任何 \(q\) ，满足 \(1 < q < a\) ）；由于 \(u_n\) 不趋于 0，有 \(\sum_{n=1}^{\infty} u_n = +\infty\) 。

这个判别法在整级数理论（我们将在以后研究它）中非常有用；但即便如此，它仍不能判定级数 \((1/n^{\alpha})\) 的收敛性，换句话说，它的应用范围比对数判别法狭窄得多。

